% Incluimos la librería que contiene las funciones auxiliares
\input{/home/julian-paz/Latex/Funciones Latex/Funciones Teoria/theoryfunctions.tex}


\usepackage{listings}
\usepackage{xcolor}

% Define colores
\definecolor{keywordcolor}{RGB}{0, 102, 204}
\definecolor{functioncolor}{RGB}{128, 0, 128}
\definecolor{constructorcolor}{RGB}{255, 69, 0}
\definecolor{typecolor}{RGB}{0, 153, 102}
\definecolor{commentcolor}{RGB}{128,128,128}
\definecolor{stringcolor}{RGB}{163,21,21}


\lstdefinelanguage{haskellcustom}{
    morekeywords={
        case, class, data, default, deriving, do, else, 
        if, import, in, infix, infixl, infixr, instance, 
        let, module, newtype, of, then, type, where
    },
    morekeywords=[2]{
        Int, Integer, Float, Double, Char, String, Bool,
        Maybe, Either, IO, Ord, Eq, Show, Read, Num,
        Functor, Applicative, Monad, Foldable, Traversable, 
        Pair, Tree, GenTree, Cont, Func, Br, Id, Expr, 
        Monoid, Output, State, T, M, StInt, N
    },
    morekeywords=[3]{
        True, False, Just, Nothing, Left, Right, 
        P, Branch, Empty, C, Gen, Funct, B, Idd, Var, 
        Add, Numb, Out, Div, Con, St, Mm, Mk, Let, State
    },
    morekeywords=[4]{
        map, filter, foldr, foldl, head, tail, take, drop,
        length, reverse, concat, zip, unzip, sum, product, 
        fmap, pure, liftA2, liftA5, SequenceA, return, foldM, 
        mapM, mempty, mappend, write, show, eval, runState, 
        runM, runC, modify, raise, doEval, get, put, throw
    },
    otherkeywords={
      ->, ::, =>, @, |, >>=
    },
    sensitive=true,
    morecomment=[l]{--},
    morecomment=[n]{\{-}{-\}},
    morestring=[b]",
    morestring=[b]´
}

\lstdefinestyle{haskellStyle}{
    language=haskellcustom,
    basicstyle=\ttfamily\small,
    keywordstyle=\color{keywordcolor}\bfseries,
    keywordstyle=[2]\color{typecolor}\bfseries,
    keywordstyle=[3]\color{constructorcolor}\bfseries,
    keywordstyle=[4]\color{functioncolor},
    commentstyle=\color{commentcolor}\itshape,
    stringstyle=\color{stringcolor},
    showstringspaces=false,
    breaklines=true
}



% Configuro el archivo de Latex
\initlatexfile{Imagenes/}

\begin{document}

\begin{flushleft}

  \begin{center}
    \textbf{\Huge{Universidad Nacional de Rosario}}
    \huge{Facultad de Ciencias Exactas, Ingenieria Y Agrimensura}\emptyline\emptyline
    \insertcenterimage{UNR}{0.8}
    \emptyline
    
    \textbf{\huge{Análisis de Lenguajes de Programación R322}}\emptyline
        
    \huge{Trabajo Práctico III - Informe}\emptyline\emptyline

    \huge{Alumnos:}
  
    \Large{Tomas Octavio Castagnino}

    \Large{Ernesto Jesús Savio}

    \Large{Julián Paz}                
  \end{center}

\newpage\initpageformat{LCC}{Análisis de Lenguajes de Programación}{Trabajo Práctico III}

\newpage\textbf{\large{Ejercicio 1. Item a}}

Probemos que la mónada State es en efecto una mónada. Para ello tenemos que probar las propiedades:

  \begin{itemize}
    \item Monad.1: Para todo $a\in A$ se cumple que:

      \begin{center}
        \begin{lstlisting}[style=haskellStyle]
            return a (>>=) k = k a
        \end{lstlisting}
      \end{center}

    \item Monad.2: 

      \begin{center}
        \begin{lstlisting}[style=haskellStyle]
            m (>>=) return = m
        \end{lstlisting}
      \end{center}

    \item Monad.3:

      \begin{center}
        \begin{lstlisting}[style=haskellStyle]
            m (>>=) (\x -> k x >>= h) = (m >>= k) >>= h
        \end{lstlisting}
      \end{center}

  \end{itemize}

Comencemos probando Monad.1, para ello consideremos $a\in A$ un elemento genérico y veamos que vale:

  \begin{center}
    $return\ a\ (>>=)\ k = k\ a$ (i)
  \end{center}

Comencemos con la prueba:\emptyline

  \hspace{5mm} $return\ a\ (>>=)\ k =_{(return.1)}$

  \hspace{5mm} $State\ (\lambda s \rightarrow (a\ :!:\ s))\ (>>=)\ k =_{(>>=.1)}$

  \hspace{5mm} $State\ (\lambda s \rightarrow \textbf{let}\ (v\ :!:\ s') = runState\ (State\ (\lambda s'' \rightarrow (a\ :!:\ s'')))\ s$
    
    \hspace{25.5mm} $\textbf{in}\ runState\ (k\ v)\ s') =_{(runState.1)}$

  \hspace{5mm} $State\ (\lambda s \rightarrow \textbf{let}\ (v\ :!:\ s') = (\lambda s'' \rightarrow (a\ :!:\ s'')))\ s$
    
    \hspace{25.5mm} $\textbf{in}\ runState\ (k\ v)\ s') =_{(\beta\ redes)}$

  \hspace{5mm} $State\ (\lambda s \rightarrow \textbf{let}\ (v\ :!:\ s') = (a\ :!:\ s)$
  
    \hspace{25.5mm} $\textbf{in}\ runState\ (k\ v)\ s') =_{(let\ in)}$

  \hspace{5mm} $State\ (\lambda s \rightarrow runState\ (k\ a)\ s) =_{(\eta\beta\ redex)}$

  \hspace{5mm} $State\ (runState\ (k\ a)) =_{(State\ \circ\ runState\ \equiv\ id)}$

  \hspace{5mm} $k\ a$\emptyline

Pasando en limpio obtenemos que:

  \begin{center}
    $return\ a\ (>>=)\ k = k\ a$ (i)
  \end{center}

Por lo tanto resulta que vale (i). Luego como $a$ es un elemento genérico resulta que para todo $a\in A$ vale la propiedad (i). 
En consecuencia vale Monad.1, que es lo que quería probar.\newpage

Ahora veamos que valga Monad.2, para ello debemos ver que se cumple que:

  \begin{center}
    $m\ (>>=)\ return = m$ (i)
  \end{center}

Comencemos con la prueba:\emptyline

  \hspace{5mm} $m\ (>>=)\ return =_{(>>=.1)}$

  \hspace{5mm} $State\ (\lambda s \rightarrow \textbf{let}\ (v\ :!:\ s') = runState\ m\ s$
  
    \hspace{25.5mm} $\textbf{in}\ runState\ (return\ v)\ s') =_{(return.1)}$

  \hspace{5mm} $State\ (\lambda s \rightarrow \textbf{let}\ (v\ :!:\ s') = runState\ m\ s$
  
    \hspace{25.5mm} $\textbf{in}\ runState\ (State\ (\lambda s'' \rightarrow (v\ :!:\ s'')))\ s') =_{(State\ \circ\ runState\ \equiv\ id)}$

  \hspace{5mm} $State\ (\lambda s \rightarrow \textbf{let}\ (v\ :!:\ s') = runState\ m\ s$
  
    \hspace{25.5mm} $\textbf{in}\ (\lambda s'' \rightarrow (v\ :!:\ s''))\ s') =_{(\beta\ redex)}$

  \hspace{5mm} $State\ (\lambda s \rightarrow \textbf{let}\ (v\ :!:\ s') = runState\ m\ s$
    
    \hspace{25.5mm} $\textbf{in}\ (v\ :!:\ s')) =_{(lema\ aux\ let-in.1)}$

  \hspace{5mm} $State\ (\lambda s \rightarrow runState\ m\ s) =_{(\eta\beta \ redex)}$

  \hspace{5mm} $State\ (runState\ m) =_{(State\ \circ\ runState\ \equiv\ id)}$
  
  \hspace{5mm} $m$\emptyline

En base a este análisis podemos concluir que:

  \begin{center}
    $m\ (>>=)\ return = m$
  \end{center}

Por lo tanto resulta que vale (i). Luego resulta que vale Monad.2, que es lo que estábamos buscando probar.\emptyline

Por último probemos que vale Monad.3, para ello tenemos que se verifica que:

  \begin{center}
    $m\ (>>=)\ (\lambda x \rightarrow k\ x\ (>>=)\ h) = (m\ (>>=)\ k)\ (>>=)\ h$ (i)
  \end{center}

Comencemos con la prueba:\emptyline

  \hspace{5mm} $m\ (>>=)\ (\lambda x\rightarrow k\ x\ (>>=)\ h) =_{(>>=.1)}$

  \hspace{5mm} $State\ (\lambda s \rightarrow \textbf{let}\ (v\ :!:\ s') = runState\ m\ s$
  
    \hspace{25.5mm} $\textbf{in}\ runState\ ((\lambda x \rightarrow k\ x\ (>>=)\ h)\ v)\ s') =_{(\beta\ redex)}$


  \hspace{5mm} $State\ (\lambda s \rightarrow \textbf{let}\ (v\ :!:\ s') = runState\ m\ s$
  
    \hspace{25.5mm} $\textbf{in}\ runState\ (k\ v\ (>>=)\ h)\ s') =_{(>>=.1)}$


  \hspace{5mm} $State\ (\lambda s \rightarrow \textbf{let}\ (v\ :!:\ s') = runState\ m\ s$
  
    \hspace{25.5mm} $\textbf{in}\ runState\ (State\ (\lambda t -> \textbf{let}\ (v'\ :!:\ s'') = runState\ (k\ v)\ t$
    
      \hspace{68.5mm} $\textbf{in}\ runState\ (h\ v')\ s''))\ s)=_{(runState\ \circ\ State\ \equiv\ id)}$


  \hspace{5mm} $State\ (\lambda s \rightarrow \textbf{let}\ (v\ :!:\ s') = runState\ m\ s$
  
    \hspace{25.5mm} $\textbf{in}\ (\lambda t -> \textbf{let}\ (v'\ :!:\ s'') = runState\ (k\ v)\ t$
    
      \hspace{42.5mm} $\textbf{in}\ runState\ (h\ v')\ s'')\ s')=_{(\beta\ redex)}$


  \hspace{5mm} $State\ (\lambda s \rightarrow \textbf{let}\ (v\ :!:\ s') = runState\ m\ s$
  
    \hspace{25.5mm} $\textbf{in}\ \textbf{let}\ (v'\ :!:\ s'') = runState\ (k\ v)\ s'$
    
      \hspace{30.5mm} $\textbf{in}\ runState\ (h\ v')\ s'')=_{(lema\ aux\ let-in.2)}$


  \hspace{5mm} $State\ (\lambda s \rightarrow \textbf{let}\ (v'\ :!:\ s'') = \textbf{let}\ (v\ :!:\ s') = runState\ m\ s$

      \hspace{52.5mm} $\textbf{in}\ runState\ (k\ v)\ s'$

    \hspace{25.5mm} $\textbf{in}\ runState\ (h\ v')\ s'')=_{(lema\ aux\ let-in.2)}$\newpage

    
  \hspace{5mm} $State\ (\lambda s \rightarrow \textbf{let}\ (v'\ :!:\ s'') = (\lambda t\rightarrow \textbf{let}\ (v\ :!:\ s') = runState\ m\ t$

      \hspace{63mm} $\textbf{in}\ runState\ (k\ v)\ s')\ s$

    \hspace{25.5mm} $\textbf{in}\ runState\ (h\ v')\ s'')=_{(runState\ \circ\ State\ \equiv\ id)}$
      

  \hspace{5mm} $State\ (\lambda s \rightarrow \textbf{let}\ (v'\ :!:\ s'') = runState\ (State\ (\lambda t\rightarrow \textbf{let}\ (v\ :!:\ s') = runState\ m\ t$

      \hspace{89mm} $\textbf{in}\ runState\ (k\ v)\ s'))\ s$

    \hspace{25.5mm} $\textbf{in}\ runState\ (h\ v')\ s'')=_{(>>=.1)}$


  \hspace{5mm} $State\ (\lambda s \rightarrow \textbf{let}\ (v'\ :!:\ s'') = runState\ (m\ (>>=)\ k)\ s$

    \hspace{25.5mm} $\textbf{in}\ runState\ (h\ v')\ s'')=_{(>>=.1)}$


  \hspace{5mm} $(m\ (>>=)\ k)\ hs$\emptyline

Pasando en limpio tenemos que:

  \begin{center}
    $m\ (>>=)\ (\lambda x \rightarrow k\ x\ (>>=)\ h) = (m\ (>>=)\ k)\ (>>=)\ h$
  \end{center}

Por lo que vale la propiedad (i), en consecuencia vale la propiedad Monad.3\emptyline

Como valen las propiedades Monad.1, Monad.2 y Monad.3, resulta que State es una mónada.\emptyline\emptyline

{\Large{\textbf{Lemas auxiliares:}}}\emptyline

  \hspace{5mm}\textbf{lema auxiliar 1 let-in:}

  \hspace{5mm} Para todo $y$ se cumple que:

  \begin{center}
    \begin{lstlisting}[style=haskellStyle]
      (let x = y in x) == y 
    \end{lstlisting}
  \end{center}

  \hspace{5mm}\textbf{lema auxiliar 2 let-in:}

  \hspace{5mm} 
  Si $a\not\in FV(h\ b)$, entonces:

  \begin{center}
    \begin{lstlisting}[style=haskellStyle]
      (let a = f x in let b = g a in h b) == (let b = let a = f x in g a in h b)
    \end{lstlisting}
  \end{center}

\newpage

\end{flushleft}
\end{document}
